% GENERATED FILE. Edit canonical lessons, then run npm run generate:books.
% canonical-source-hash: fnv1a64:3e10027fa2a88b20
% canonical-lessons: TE-C307-phrijju, TE-C307-miksi, TE-C307-antlu, TE-C307-agarubatti, TE-C307-kunkuma

\chapter{Fridge, Mixer-grinder, Dirty dishes, Incense stick, Kumkum}
\label{ch:te-fridge-mixer-grinder-dirty-dishes-incense-stick-kumkum}
\hlchaptermodality{307}

\begin{chapteropening}
\textbf{By the end of this chapter:} \emph{I can say a fridge, a mixer-grinder, dirty dishes, an incense stick and kumkum.}

Complete the last lesson of chapter 307: I can say a fridge, a mixer-grinder, dirty dishes, an incense stick and kumkum.
\end{chapteropening}

\section[\texorpdfstring{\te{ఫ్రిజ్జు} (phrijju)}{ఫ్రిజ్జు (phrijju)}]{\te{ఫ్రిజ్జు} (phrijju) — a fridge}

\label{lesson:TE-C307-phrijju}



\begin{quote}
Before the new one: say the Telugu for cotton, then the Telugu for wool.
\end{quote}

\subsection*{You'll want to know: \te{ఫ్రిజ్జు}}
\textbf{\te{ఫ్రిజ్జు}} — \emph{phrijju} — \textquotedblleft{}a fridge\textquotedblright{}.

\begin{grammarlens}[title={What you've built}]
One more word for this chapter.
\end{grammarlens}

\subsection*{Your turn}
\begin{itemize}
\raggedright
  \item \emph{Say it:} \emph{phrijju}
  \item \emph{Say it:} \emph{phrijju}, once more
  \item \emph{Say it:} say \emph{phrijju}
  \item \emph{Recall:} say the Telugu for cotton, then the Telugu for wool, then say \emph{phrijju} again
\end{itemize}

\begin{tcolorbox}[breakable,colback=teal!4,colframe=teal!35!black,title={Before you move on}]
What does \te{ఫ్రిజ్జు} mean? (\textquotedblleft{}A fridge\textquotedblright{}.) Say it once more.
\end{tcolorbox}
\section[\texorpdfstring{\te{మిక్సీ} (miksī)}{మిక్సీ (miksī)}]{\te{మిక్సీ} (miksī) — a mixer-grinder}

\label{lesson:TE-C307-miksi}



\begin{quote}
Before the new one: say the Telugu for wool, then the Telugu for a fridge.
\end{quote}

\subsection*{You'll want to know: \te{మిక్సీ}}
\textbf{\te{మిక్సీ}} — \emph{miksī} — \textquotedblleft{}a mixer-grinder\textquotedblright{}.

\begin{grammarlens}[title={What you've built}]
One more word for this chapter.
\end{grammarlens}

\subsection*{Your turn}
\begin{itemize}
\raggedright
  \item \emph{Say it:} \emph{miksī}
  \item \emph{Say it:} \emph{miksī}, once more
  \item \emph{Say it:} say \emph{miksī}
  \item \emph{Recall:} say the Telugu for wool, then the Telugu for a fridge, then say \emph{miksī} again
\end{itemize}

\begin{tcolorbox}[breakable,colback=teal!4,colframe=teal!35!black,title={Before you move on}]
What does \te{మిక్సీ} mean? (\textquotedblleft{}A mixer-grinder\textquotedblright{}.) Say it once more.
\end{tcolorbox}
\section[\texorpdfstring{\te{అంట్లు} (aṇṭlu)}{అంట్లు (aṇṭlu)}]{\te{అంట్లు} (aṇṭlu) — dirty dishes}

\label{lesson:TE-C307-antlu}



\begin{quote}
Before the new one: say the Telugu for a fridge, then the Telugu for a mixer-grinder.
\end{quote}

\subsection*{You'll want to know: \te{అంట్లు}}
\textbf{\te{అంట్లు}} — \emph{aṇṭlu} — \textquotedblleft{}dirty dishes\textquotedblright{}.

\begin{grammarlens}[title={What you've built}]
One more word for this chapter.
\end{grammarlens}

\subsection*{Your turn}
\begin{itemize}
\raggedright
  \item \emph{Say it:} \emph{aṇṭlu}
  \item \emph{Say it:} \emph{aṇṭlu}, once more
  \item \emph{Say it:} say \emph{aṇṭlu}
  \item \emph{Recall:} say the Telugu for a fridge, then the Telugu for a mixer-grinder, then say \emph{aṇṭlu} again
\end{itemize}

\begin{tcolorbox}[breakable,colback=teal!4,colframe=teal!35!black,title={Before you move on}]
What does \te{అంట్లు} mean? (\textquotedblleft{}Dirty dishes\textquotedblright{}.) Say it once more.
\end{tcolorbox}
\section[\texorpdfstring{\te{అగరుబత్తి} (agarubatti)}{అగరుబత్తి (agarubatti)}]{\te{అగరుబత్తి} (agarubatti) — an incense stick}

\label{lesson:TE-C307-agarubatti}



\begin{quote}
Before the new one: say the Telugu for a mixer-grinder, then the Telugu for dirty dishes.
\end{quote}

\subsection*{You'll want to know: \te{అగరుబత్తి}}
\textbf{\te{అగరుబత్తి}} — \emph{agarubatti} — \textquotedblleft{}an incense stick\textquotedblright{}.

\begin{grammarlens}[title={What you've built}]
One more word for this chapter.
\end{grammarlens}

\subsection*{Your turn}
\begin{itemize}
\raggedright
  \item \emph{Say it:} \emph{agarubatti}
  \item \emph{Say it:} \emph{agarubatti}, once more
  \item \emph{Say it:} say \emph{agarubatti}
  \item \emph{Recall:} say the Telugu for a mixer-grinder, then the Telugu for dirty dishes, then say \emph{agarubatti} again
\end{itemize}

\begin{tcolorbox}[breakable,colback=teal!4,colframe=teal!35!black,title={Before you move on}]
What does \te{అగరుబత్తి} mean? (\textquotedblleft{}An incense stick\textquotedblright{}.) Say it once more.
\end{tcolorbox}
\section[\texorpdfstring{\te{కుంకుమ} (kuṅkuma)}{కుంకుమ (kuṅkuma)}]{\te{కుంకుమ} (kuṅkuma) — kumkum, vermilion powder}

\label{lesson:TE-C307-kunkuma}



\begin{quote}
Before the new one: say the Telugu for dirty dishes, then the Telugu for an incense stick.
\end{quote}

\subsection*{You'll want to know: \te{కుంకుమ}}
\textbf{\te{కుంకుమ}} — \emph{kuṅkuma} — \textquotedblleft{}kumkum, vermilion powder\textquotedblright{}.

\begin{grammarlens}[title={What you've built}]
One more word for this chapter.
\end{grammarlens}

\subsection*{Your turn}
\begin{itemize}
\raggedright
  \item \emph{Say it:} \emph{kuṅkuma}
  \item \emph{Say it:} \emph{kuṅkuma}, once more
  \item \emph{Say it:} say \emph{kuṅkuma}
  \item \emph{Recall:} say the Telugu for dirty dishes, then the Telugu for an incense stick, then say \emph{kuṅkuma} again
\end{itemize}

\begin{tcolorbox}[breakable,colback=teal!4,colframe=teal!35!black,title={Before you move on}]
What does \te{కుంకుమ} mean? (\textquotedblleft{}Kumkum, vermilion powder\textquotedblright{}.) Say it once more.
\end{tcolorbox}
